\chapter{Gaussian kernels and stationary backgrounds}
\label{ch05:gaussian}

The quadratic approximation serves two purposes. It identifies the
covariance of the fluctuations eliminated by a block transformation, and
it supplies the inverse needed to construct a stationary background for
the nonlinear action. These uses require different estimates. Exact matrix
identities hold on each finite lattice; uniform locality requires bounds
which survive changes of the lattice size and of the blocking history.

We first fix the normalization of the exact identities. We then isolate the
additional hypotheses for uniform analytic continuation and for a
stationary background. The general block transformation algebra is
standard; see Ba{\l}aban--Feldman--Kn\"orrer--Trubowitz
\cite{BalabanFeldmanKnoerrerTrubowitz2016}. The construction of local
propagators and backgrounds for lattice gauge theory originates in
Ba{\l}aban's work \cite{Balaban1984Propagators,Balaban1985Variational}.
Those references explain the method. Their results cannot be applied to a
different averaging map without checking its hypotheses.

\section{Lattice differential forms and zero modes}
\label{ch05:forms}

In this chapter distances and Fourier variables are in lattice units.
Let $\Lambda_s=(\ZZ/s\ZZ)^4$. A lattice $r$-form assigns an element of
$\gLie$ to each site and each increasing $r$-tuple of coordinate
directions. The scalar product is the sum over sites and increasing
tuples, using the Lie algebra metric of \eqref{ch01:metric}.
In particular there is no volume normalization in this scalar product.
Write $\Omega^r(\Lambda_s;\gLie)$ for this finite-dimensional space.
The forward difference is
$\nabla_\mu u(x)=u(x+e_\mu)-u(x)$, and the exterior derivative is
\[
 (d_r u)_{\mu_0\cdots\mu_r}
 =\sum_{i=0}^r(-1)^i\nabla_{\mu_i}
       u_{\mu_0\cdots\widehat{\mu_i}\cdots\mu_r}.
\]
Commutation of the differences gives $d_{r+1}d_r=0$.
An adjoint is always taken for the stated counting scalar products.

A linearized gauge transformation of a one-form $a$ is
$a\mapsto a+d_0\phi$. If $H$ is a positive self-adjoint operator on
two-forms, the quadratic action and its kinetic operator are
\begin{equation}
 S_0(a)=\tfrac12\langle d_1a,Hd_1a\rangle,\qquad K=d_1^*Hd_1.
 \label{ch05:kinetic}
\end{equation}
Thus $K$ vanishes on gradients. It also vanishes on constant one-forms.
On a periodic lattice the latter are harmonic modes and are not gradients
of periodic scalar functions. Removing gauge directions alone therefore
does not make $K$ invertible.

For $p\in(2\pi/s)\ZZ^4/(2\pi\ZZ)^4$ our Fourier convention is
\begin{equation}
 \widehat u(p)=\sum_{x\in\Lambda_s}e^{-ip\cdot x}u(x),
 \quad
 u(x)=s^{-4}\sum_p e^{ip\cdot x}\widehat u(p),
 \quad
 \|u\|_2^2=s^{-4}\sum_p|\widehat u(p)|^2.
 \label{ch05:fourier}
\end{equation}
Set $q_\mu(p)=e^{ip_\mu}-1$ and
$\lambda(p)=\sum_\mu q_\mu(-p)q_\mu(p)$.
On real momenta this is $\sum_\mu|q_\mu(p)|^2$.
The symbol of $d_r$ is exterior multiplication by $q$.
In particular, on one-forms,
\[
 |q\wedge a|^2+|\overline q\mathbin{\lrcorner}a|^2
      =\lambda(p)|a|^2.
\]
Here contraction denotes the adjoint of exterior multiplication.
For $p\ne0$ this gives the orthogonal decomposition into a gradient
and its orthogonal complement. At $p=0$ the whole one-form space is
harmonic.

Whenever a Gaussian probability measure is used below, its precision is
strictly positive on the specified coordinate space. One may work on a
fixed slice with the harmonic modes removed, or add a positive regulator.
A massless covariance at zero momentum is not defined by writing
$1/\lambda(0)$. Conversely, an identity between quadratic forms can
remain valid on a space where no normalized massless Gaussian exists.

\section{The exact Gaussian block transformation}
\label{ch05:schur}

Let $E_f,E_c$ be real finite-dimensional inner-product spaces, $Q:E_f\to
E_c$ a linear block map, and $b>0$ the block penalty. Suppose initially
that $K_f>0$ on $E_f$. Define
\begin{equation}
 A=K_f+bQ^*Q,\qquad
 W=A^{-1}bQ^*,\qquad
 K_c=bI-b^2QA^{-1}Q^* .
 \label{ch05:schur-data}
\end{equation}
These formulas describe the normalized Gaussian problem. When the
physical quadratic exponent is
$\beta(\langle A_f,K_fA_f\rangle+b\|QA_f-V\|^2)/2$,
the variables here are $a=\sqrt{\beta}A_f$ and
$v=\sqrt{\beta}V$. In unscaled link logarithms the fluctuation
covariance is $\beta^{-1}(K_f+bQ^*Q)^{-1}$.
Thus the block parameter remains $b$ after normalization, although
the original block kernel has penalty $b\beta$.

The map $W$ sends a coarse field to the stationary fine field. The
fluctuation covariance is $A^{-1}$. The Schur-complement formulas give
\begin{align}
 K_c&=(b^{-1}I+QK_f^{-1}Q^*)^{-1}, \label{ch05:coarse-cov}\\
 \langle a,K_fa\rangle+b\|Qa-v\|^2
 &=\langle v,K_cv\rangle+\langle a-Wv,A(a-Wv)\rangle .
 \label{ch05:completion}
\end{align}
These are the usual finite-dimensional Gaussian identities, or the
real quadratic specialization of the block transformation identities in
\cite{BalabanFeldmanKnoerrerTrubowitz2016}; no separate proof is needed.
The accompanying normalization is
\[
 \int_{E_f}e^{-\langle a,K_fa\rangle/2-b\|Qa-v\|^2/2}\,da
 =(2\pi)^{\dim E_f/2}(\det A)^{-1/2}
            e^{-\langle v,K_cv\rangle/2}.
\]
The Lebesgue measure here comes from the specified inner product.
In a normalized block kernel there is also the factor
$(b/2\pi)^{\dim E_c/2}$.

There is an equivalent probabilistic statement. If $a$ and $\xi$ are
independent centered Gaussians with covariances $K_f^{-1}$ and $I$,
then $v=Qa+b^{-1/2}\xi$ has covariance
$b^{-1}I+QK_f^{-1}Q^*$. Conditional on $v$, the fine field has mean
$Wv$ and covariance $A^{-1}$.

The identities \eqref{ch05:schur-data} and
\eqref{ch05:completion} also hold for $K_f\ge0$, provided $A$ is
invertible. Their derivation uses no inverse of $K_f$.
In this case \eqref{ch05:coarse-cov} and the joint probability
interpretation are unavailable unless the remaining null directions have
been treated. In the gauge problem, every use of this extension must
specify the slice or quotient on which $A$ is invertible.

The useful consequences, with the same domains and adjoints, are
\[
 K_fW=Q^*K_c,\qquad QW=I-b^{-1}K_c,\qquad
 W^*K_fW+b(QW-I)^*(QW-I)=K_c.
\]
The last identity says that the stationary fine action together with
its block penalty is exactly the coarse quadratic action.

\section{Averaging and the Fourier alias formula}
\label{ch05:alias}

To check the scaling factor in the covariance recursion, take
$\Lambda_{Ls}$ as fine lattice and $\Lambda_s$ as coarse lattice, with
odd $L\ge3$. Let $w(z)$ be a finitely supported real averaging kernel
of total mass one, periodized on the fine lattice, and use
\[
 (Ma)_\mu(y)=\sum_z w(z)\sum_{r=0}^{L-1}
                    a_\mu(Ly+z+re_\mu).
\]
The inner sum is a sum along a link of length $L$, not its average.
Define the symbols
\[
 \widehat w(\theta)=\sum_z w(z)e^{i\theta\cdot z},
 \qquad r_L(t)=L^{-1}\sum_{r=0}^{L-1}e^{irt}.
\]
The sign in $\widehat w$ is chosen for this formula for $M$.
The aliases of a coarse momentum $p$ are
$\theta_z=(p+2\pi z)/L$, with one representative for each
$z\in(\ZZ/L\ZZ)^4$.

Suppose an auxiliary, translation-invariant Gaussian one-form has
diagonal covariance symbol $\operatorname{diag}(c_\mu(\theta))$ in
the direction index. Color indices are suppressed and are unchanged
by this calculation. The symbol of the covariance of
$Ma+b^{-1/2}\xi$ is
\begin{equation}
 c'_\mu(p)=b^{-1}
 +L^{-2}\sum_{z\ {\rm mod}\ L}
 \widehat w(\theta_z)\widehat w(-\theta_z)
 r_L(\theta_{z,\mu})r_L(-\theta_{z,\mu})c_\mu(\theta_z).
 \label{ch05:alias-recursion}
\end{equation}
This formula is exact for a positive regulated covariance. It can
also be used for nonzero modes of a massless symbol with the zero
mode handled separately.

For clarity, the normalization check is short. Substitution of
\eqref{ch05:fourier} into $Ma$ gives
$\widehat{Ma}_\mu(p)=L^{-4}\sum_z
L\widehat w(\theta_z)r_L(\theta_{z,\mu})
\widehat a_\mu(\theta_z)$.
Fine Fourier covariances carry the volume $(Ls)^4$, whereas coarse
ones carry $s^4$. Consequently the covariance factor is
$L^{-8}L^2L^4=L^{-2}$. This is the factor used in
\eqref{ch05:alias-recursion}.

The particular smoothing in the supplied construction uses
\[
 \widehat w_L(\theta)=
 \prod_{\nu=1}^4\left(L^{-1}
       \sum_{|r|\le(L-1)/2}e^{ir\theta_\nu}\right),
 \qquad \widehat w(\theta)=\widehat w_L(\theta)^j,
\]
where $j$ is the integer smoothing order. Centering makes
$\widehat w_L$ even on real and complex momenta.
The recurrence then contains $\widehat w_L(\theta_z)^{2j}$.
This choice fixes an auxiliary covariance. Identification with the
quadratic form of a particular gauge-equivariant block map additionally
requires checking that its linearization differs from $M$ only by an
appropriate gradient, and that the chosen treatment of zero modes agrees.

\section{Gauge reduction and a regular curvature symbol}
\label{ch05:curvature}

At a real nonzero momentum suppose $c_\mu>0$ and put
$w_\mu=c_\mu^{-1}$. Minimizing over a scalar gauge parameter gives
\begin{equation}
 \min_\phi\sum_\mu w_\mu|a_\mu+q_\mu\phi|^2
 =\frac{\displaystyle\sum_{\mu<\nu}w_\mu w_\nu
                  |q_\mu a_\nu-q_\nu a_\mu|^2}
        {\displaystyle\sum_\mu w_\mu|q_\mu|^2}.
 \label{ch05:quotient}
\end{equation}
It is the weighted orthogonal projection formula. Thus a diagonal
two-form metric representing the quotient form is
\begin{equation}
 (H_0)_{\mu\nu}(p)=
 \left[c_\mu(p)c_\nu(p)
               \sum_\alpha\frac{q_\alpha(-p)q_\alpha(p)}
                                  {c_\alpha(p)}\right]^{-1}.
 \label{ch05:H0}
\end{equation}
Formula \eqref{ch05:H0} is initially an identity away from zero.
Its regular extension at zero is a separate issue.

Here is the cancellation used for that extension. Let $M$ now denote
the total dilation over several steps; it is a number, distinct from
the averaging operator of the previous section.
Put $\lambda_M(p)=M^2\lambda(p/M)$.
Suppose the accumulated covariance has the representation
\[
 c_\mu(p)=\frac{A_\mu(p)}{\lambda_M(p)}+G_\mu(p),
 \quad
 A_\mu(p)=F(p)r_M(p_\mu/M)r_M(-p_\mu/M),
 \quad B_\mu=A_\mu+\lambda_MG_\mu.
\]
The remainder $G_\mu$ collects the nonprincipal aliases and the
accumulated block noise. Assume that $F,G_\mu,A_\mu,B_\mu$ are
holomorphic on a common neighborhood and that $F,A_\mu,B_\mu$ do
not vanish there. The elementary identity
$q_\mu(p)r_M(p_\mu/M)^{-1}=M q_\mu(p/M)$, interpreted after
multiplication to remove possible zeros, gives
\[
 q_\mu(p)q_\mu(-p)F
       =M^2q_\mu(p/M)q_\mu(-p/M)A_\mu .
\]
Consequently, if
\begin{equation}
 h=F^{-1}-\sum_\mu
       \frac{q_\mu(p)q_\mu(-p)G_\mu}{A_\mu B_\mu},
 \qquad
 \lambda_M h=\sum_\mu\frac{q_\mu(p)q_\mu(-p)}{B_\mu},
 \label{ch05:pole-cancel}
\end{equation}
then \eqref{ch05:H0} becomes
\[
                 (H_0)_{\mu\nu}=(h B_\mu B_\nu)^{-1}.
\]
The second equality in \eqref{ch05:pole-cancel} follows by subtracting
$1/B_\mu$ from $1/A_\mu$; it is not an additional assumption.
If $h$ is also nonvanishing, the last formula proves holomorphic
extension through the apparent massless pole.
When $F(0)=A_\mu(0)=1$, it gives $H_0(0)=I$.

This calculation proves pole cancellation, not a uniform zero-free
strip. For that stronger conclusion the nonprincipal aliases,
accumulated noise, and reciprocal factors must be bounded uniformly
over every admitted blocking history. These are precisely the
estimates which cannot be replaced by finite-dimensional Gaussian
algebra.

One can modify the metric away from exact two-forms without changing
the action. For $\zeta\ge0$ define, at real momentum,
\begin{equation}
 H^{-1}=H_0^{-1}
       +\zeta(d_2H_0^{-1})^*(d_2H_0^{-1}).
 \label{ch05:metric-modification}
\end{equation}
If $H_0>0$, this defines a positive $H$, and
$H^{-1}H_0d_1=d_1$, because $d_2d_1=0$. Therefore
$Hd_1=H_0d_1$. Positivity on all two-forms and agreement on exact
two-forms are different assertions; this modification is useful
because later estimates need both.

\section{Uniform estimates over finite blocking histories}
\label{ch05:histories}

We now verify the uniform reciprocal bounds required in
\eqref{ch05:pole-cancel} for the actual averaging recurrence.
A finite history consists of odd ratios $L_1,\ldots,L_n\ge3$ and
smoothing orders $j_0,\ldots,j_n\ge4$, with
$j_{h-1}-j_h\in\{0,1\}$. Set $M_h=L_1\cdots L_h$, $M_0=1$.
The order used by the $h$th block map is its incoming order
$j_{h-1}$. Denote the centered box symbol $\widehat w_L$ of
Section~\ref{ch05:alias} by $w_L$ in the following formulas.

\begin{proposition}[Uniform curvature metric]
\label{ch05:history-prop}
Fix the final smoothing order $j=j_n\ge4$ and the penalty $b>0$.
Start with the nonzero-mode covariance $c_\mu(p)=1/\lambda(p)$
and iterate \eqref{ch05:alias-recursion} along any such finite
history. The resulting curvature metric $H_0$ has a holomorphic,
periodic extension to a common strip about real momenta, with
uniformly bounded entries and inverse entries. On real momenta
$H_0$ has uniform positive upper and lower bounds and $H_0(0)=I$.
The constants depend on $j,b$, but not on the history length,
ratios or torus period.
\end{proposition}

\begin{proof}
The box composition identity
$w_A(B\theta)w_B(\theta)=w_{AB}(\theta)$ holds for odd
$A,B$, by decomposing a centered interval of length $AB$
into centered intervals of length $B$.
Thus the accumulated filter is
\begin{equation}
 f(\theta)=w_{M_n}(\theta)^j
       \prod_{\{h:j_{h-1}=j_h+1\}}w_{M_h}(\theta).
 \label{ch05:history-filter}
\end{equation}
The directional filter is
$f(\theta)r_{M_n}(\theta_\mu)$: the link factors compose by
$r_A(Bt)r_B(t)=r_{AB}(t)$.

Here and below momentum estimates are on an enlarged cube
$|\operatorname{Re}p_\mu|<\pi+\eta$,
$|\operatorname{Im}p_\mu|<\eta$, for a sufficiently small fixed
$\eta>0$. Write $M=M_n$ and choose centered alias representatives
$z_\mu\in[-(M-1)/2,(M-1)/2]$. The finite geometric sum gives
\[
 |r_M((p_\mu+2\pi z_\mu)/M)|,\quad
 \left|M^{-1}\sum_{|r|\le(M-1)/2}
     e^{ir(p_\mu+2\pi z_\mu)/M}\right|
       \le \frac{C}{1+|z_\mu|}.
\]
Indeed the numerator of the geometric quotient is bounded, and
$M|e^{i(p_\mu+2\pi z_\mu)/M}-1|$ is bounded below by
$c(1+|z_\mu|)$ for $z_\mu\ne0$; the principal case is bounded
directly by the finite sum. At the principal real alias the
centered factor has a positive lower bound.

The additional factors in \eqref{ch05:history-filter} have a
uniform product bound. A factor of side $A\le M$ has modulus
at most $\exp(C\eta A/M)$ on the cube evaluated at $p/M$
or any other alias. The sides occurring in the product increase
at least geometrically, so their sum divided by $M$ is bounded.
On the real principal alias the same factors have product bounded
away from zero. To see this without multiplying a fixed lower
bound infinitely often, use
$|\log w_A(p/M)|\le C A/M$ on that real cube, with the value
at a removable zero interpreted continuously. The bound follows
from the geometric quotient near zero and compactness for
$A/M$ bounded below. Again $\sum A/M$ is bounded.
Consequently, with
$\Gamma(z)=\prod_\mu(1+|z_\mu|)$,
\[
 |f((p+2\pi z)/M)r_M((p_\mu+2\pi z_\mu)/M)|
       \le C_j\Gamma(z)^{-j},
 \qquad |f(p/M)|\ge c_j>0
\]
in the first estimate on a complex cube and in the second on its
real slice.

There is also an exact description of the accumulated noise.
The independent noise introduced at step $h$ is propagated only
through the suffix $h+1,\ldots,n$. If that suffix has total ratio
$D=M_n/M_h$, its contribution, before multiplication by $b^{-1}$,
is
\[
 D^{-2}\sum_{z\ {\rm mod}\ D}
   f_{\rm suffix}(\theta_z)f_{\rm suffix}(-\theta_z)
   r_D(\theta_{z,\mu})r_D(-\theta_{z,\mu}).
\]
This includes $1$ when $h=n$ and $D=1$.
To verify the formula, apply \eqref{ch05:alias-recursion} to one
constant covariance and split each residue modulo a product of
ratios into its successive residues. The filter composition
identities above identify the summands and the factors multiply
to $D^{-2}$. Summing over the noises gives precisely the noise
part of the iterated recurrence.
Each suffix contribution is bounded by $C_jD^{-2}$ because
\[
 \sum_z\Gamma(z)^{-2j}\le
 \left(\sum_{k\in\ZZ}(1+|k|)^{-2}\right)^4<81.
\]
As $D\ge3^{n-h}$, the total noise is at most
$C_j\sum_{r\ge0}3^{-2r}$, uniformly in $n$.

Separate the principal alias of the propagated initial
covariance from all other aliases. The former gives
$A_\mu/\lambda_M$ with the notation of
\eqref{ch05:pole-cancel}; the latter are
\[
 \sum_{z\ne0}
 \frac{f(\theta_z)f(-\theta_z)
                 r_M(\theta_{z,\mu})r_M(-\theta_{z,\mu})}
                      {M^2\lambda(\theta_z)}.
\]
For a nonprincipal real alias the denominator is bounded below
by $c(1+\max_\mu|z_\mu|)^2$. Its change on a complex strip of
width $\eta$ is at most
$C\eta(1+\max_\mu|z_\mu|)+C\eta^2$.
Decreasing $\eta$ gives the same lower bound in modulus.
The preceding summability therefore bounds this sum holomorphically.
Together with the suffix noise estimate it proves a uniform
holomorphic bound for $G_\mu$, which is nonnegative on the real
slice.

Recall that $F=f(p/M)f(-p/M)$,
$A_\mu=F r_M(p_\mu/M)r_M(-p_\mu/M)$ and
$B_\mu=A_\mu+\lambda_MG_\mu$.
The principal filter bounds give positive real lower bounds for
$F,A_\mu$, and hence for $B_\mu$. All three have uniform complex
upper bounds. Cauchy estimates on a slightly larger cube therefore
give uniformly bounded reciprocals on a smaller one.
The identity in \eqref{ch05:pole-cancel} now supplies the real
lower bound for $h$: since $B_\mu\le C_B$,
\[
 h(p)=\frac{\sum_\mu |q_\mu(p)|^2/B_\mu(p)}{\lambda_M(p)}
          \ge c/C_B\quad(p\ne0).
\]
We used uniform comparison of both $\lambda(p)$ and
$\lambda_M(p)$ with $|p|^2$ on the real cube. At zero $h=1$.
Its defining formula gives a holomorphic upper bound using the
reciprocals already obtained, and another Cauchy estimate gives
a uniformly bounded reciprocal for $h$.
Thus $(H_0)_{\mu\nu}=(hB_\mu B_\nu)^{-1}$ has all the asserted
bounds. The order of this argument is important: the lower bound
for $h$ uses $B_\mu$, not an assumed inverse bound for $h$.
Periodicity holds first on real overlaps by permutation of the
aliases in the original covariance recurrence. Analytic continuation
then joins the cubes to a strip. The same Fourier symbol works on
every finite-period grid.
\end{proof}

This proves uniform bounds for arbitrary finite histories.
Convergence for a prescribed sequence of growing histories is
stronger and is not a consequence of this proposition alone.

\section{Exact transports and the curvature defect}
\label{ch05:transports}

Recall that $H_0$ is the two-form metric in \eqref{ch05:H0}; $H$ denotes its positive completion from \eqref{ch05:metric-modification}, agreeing with $H_0$ on exact two-forms. Fine and coarse subscripts below specify the scale of each metric.

We describe the algebra at one real nonzero coarse momentum.
This also specifies which transport must be estimated. The use
of aliases and block-periodic Fourier coordinates is standard;
see \cite{BalabanFeldmanKnoerrerTrubowitzBloch2016}.
Write $\theta_l=(p+2\pi l)/L$,
$\delta_r=Ld_r(\theta_l)$ for the rescaled fine exterior derivative,
and $d_r=d_r(p)$ for the coarse one.
Let
\[
 D_l=L^{-2}\operatorname{diag}(c_{f,\mu}(\theta_l)),
 \quad D=\operatorname{diag}(c_{c,\mu}(p)),\quad
 x=d_0,\quad x_l=\delta_0 .
\]
For this paragraph $H,H_l$ denote the unmodified curvature metrics.
Set
\[
 R=d_1^*Hd_1,\quad R_l=\delta_1^*H_l\delta_1,\quad
 \rho=x^*D^{-1}x,\quad \pi=\rho^{-1}x^*D^{-1}.
\]
The gauge quotient identity gives
\[
 R=D^{-1}-\rho^{-1}D^{-1}xx^*D^{-1},\qquad
 DR=I-x\pi .
\]
The fine formulas have subscript $l$. Define
\[
 m_l=w_L(\theta_l)^j,\qquad
 F_{r,l}=m_l\bigwedge\nolimits^r
            \operatorname{diag}(r_L(\theta_{l,\mu})).
\]
The link identity and the covariance recurrence say, respectively,
\begin{equation}
 F_{r+1,l}\delta_r=d_rF_{r,l},\qquad
             \sum_lF_{1,l}D_lF_{1,l}^*=D-b^{-1}I .
 \label{ch05:transport-input}
\end{equation}
Introduce the positive scalar
$\alpha=\sum_l|m_l|^2\tr D_l$, and put
$e_l=\overline m_l\tr D_l/\alpha$,
$y_l=|m_l|^2\rho/\rho_l$, and $C=Hd_1Dd_1^*$.
The rescaled one-form and two-form transports are
\begin{align}
 \mathcal A_l&=D_lF_{1,l}^*R+e_lx_l\pi,\\
 \mathcal T_l&=H_l\delta_1D_lF_{1,l}^*d_1^*
                         +y_lF_{2,l}^*(I-C).
 \label{ch05:transport-formula}
\end{align}
Their defining identities are
\begin{align}
 R_l\mathcal A_l&=F_{1,l}^*R,&
 \delta_1^*\mathcal T_l&=F_{1,l}^*d_1^*,&
 \mathcal T_lHd_1&=H_l\delta_1\mathcal A_l,\\
 \sum_lF_{1,l}\mathcal A_l&=I-b^{-1}R.
 \label{ch05:transport-identities}
\end{align}
These follow by substituting the quotient formula for $R$ and
using \eqref{ch05:transport-input}; in the last equality the
gradient contribution is $x\pi\sum_lm_le_l=x\pi$.
In particular the term $I-C$ cannot simply be omitted: it is
annihilated by $d_1^*$ and annihilates $Hd_1$, but controls the
transport on the complementary two-forms.

The physical kernels from a coarse site to a fine site $Ly+r$
are obtained from the alias sums
\[
 L^{-1}\sum_l e^{i\theta_l\cdot r}\mathcal A_l(p),
 \qquad
 L^{-2}\sum_l e^{i\theta_l\cdot r}\mathcal T_l(p)
\]
by coarse Fourier inversion. Denote the resulting operators by
$W$ and $T$. Restoring counting adjoints gives
\begin{equation}
 T^*d_{1,f}=d_{1,c}Q,\qquad
 TH_cd_{1,c}=H_fd_{1,f}W,\qquad
 QW=I-b^{-1}d_{1,c}^*H_cd_{1,c}.
 \label{ch05:physical-transports}
\end{equation}
Here $Q$ is the straight averaging operator in this construction.
The same statements for a block average containing a gradient
correction require its corresponding choice of the gradient
part of $W$.

The apparent poles in \eqref{ch05:transport-formula} can be
removed using the principal covariance split. We record the
crucial cancellation to avoid concealing a uniformity assumption.
Write $s=\lambda_{M_c}(p)$, $D=s^{-1}\operatorname{diag}(B_\mu)$,
$D_0=s^{-1}\operatorname{diag}(B_{0,\mu})$, and
\[
 D=F_{1,0}D_0F_{1,0}^*+U,\qquad U=\operatorname{diag}(u_\mu).
\]
The tail $U$ consists of nonprincipal aliases and the current
noise; it is uniformly holomorphic by the estimates just proved.
Let $z_\mu=u_\mu/B_\mu$. Then
\[
 \frac{|m_0r_L(\theta_{0,\mu})|^2B_{0,\mu}}{B_\mu}
                      =1-sz_\mu .
\]
With $J=F_{2,0}^{-1}$, define the bounded analytic matrix
\[
 \mathcal B=H_0\left[
       \operatorname{diag}_{\mu<\nu}(z_\mu+z_\nu)
                    -s\bigwedge\nolimits^2\operatorname{diag}(z_\mu)
                         \right].
\]
Here $H_0$ means the fine metric at alias zero, not a new
modification of the metric. Entrywise multiplication gives
$s\mathcal B=H_0-y_0F_{2,0}^*HF_{2,0}$, and hence
\[
 \mathcal T_0=y_0F_{2,0}^*
   +\bigl(\mathcal B Jd_1\operatorname{diag}(B_\mu)
                                -H_0Jd_1U\bigr)d_1^*.
\]
All factors are regular at zero. For $\mathcal A_0$, the
corresponding principal denominator is
$s\alpha=|m_0|^2\sum_\mu B_{0,\mu}+s\alpha_{\ne0}$,
which is bounded below on the real cube. Its reciprocal is
therefore uniformly holomorphic. Subtracting the leading
diagonal term in $\mathcal A_0$ leaves the regular factor
\[
 m_0^{-1}\left(u_\mu/B_\mu-\alpha_{\ne0}/(s\alpha)\right).
\]
For nonprincipal aliases the denominators have the lower bounds
used in Proposition~\ref{ch05:history-prop}.
These formulas give bounded principal columns and bounds
$C\Gamma(l)^{-j}$ for both nonprincipal columns on a smaller cube.
Adjoints in these analytic formulas mean transpose at opposite
momentum, with $l$ replaced by $-l$.

Each fine difference in the alias sum contributes
$e^{i\theta_{l,\mu}}-1$, bounded by
$CL^{-1}(1+\max_\nu|l_\nu|)$.
The summability estimate for $\Gamma$ therefore gives, for
$r\le j-2$, exponentially weighted coarse Fourier coefficient
bounds of order $L^{-1-r}$ for $\nabla_f^rW$ and $L^{-2-r}$
for $\nabla_f^rT$. Fourier contour displacement is the usual
conversion of strip analyticity to exponential coefficient
decay. It is performed after summing the aliases, whose
permutation under $p\mapsto p+2\pi n$ gives periodicity.
When a fine difference crosses a block boundary, its coarse
argument is shifted by one block; this is included by the phase
$e^{i\theta_l\cdot r}$ and must not be dropped.
At zero momentum the principal columns are $I$ and the other
columns vanish, so the identities extend there without choosing
a massless Gaussian normalization.

\subsection{A positive factorization of the defect}

The next argument uses exterior-complex exactness together with
the bounds just obtained. Exactness for a regular sequence is
standard Koszul algebra \cite[Lemma 15.31.2]{StacksKoszul}.
The quantitative version needed here preserves exponential
coefficient bounds: if an exponentially summable matrix row
$A$ on positive-degree forms satisfies $Ad_r=0$, it can be
written $A=Vd_{r+1}$ on a smaller strip, with a fixed norm cost.
One obtains this version by successively dividing
$A(p)-A(p)|_{p_a=0}$ by $e^{ip_a}-1$ and contracting the
$a$th exterior index. The elementary quotient
\[
 \frac{e^{inp_a}-1}{e^{ip_a}-1}
   =\sum_{h=0}^{n-1}e^{ihp_a}\quad(n>0)
\]
and the corresponding negative sum for $n<0$ show that a
division costs at most $C/\eta$ when width $\eta$ is reduced
to $\eta/2$. There are only four coordinates. This is the
usual coordinate proof of Koszul exactness, with the coefficient
norm retained. Applying it on both sides shows that a symmetric
two-form row $E$ with $Ed_1=d_1^*E=0$ has
$E=d_2^*C_0d_2$, with $C_0$ symmetric and exponentially summable.
For completeness, first choose a left lift $d_2^*W_0=E$.
Since $d_2^*W_0d_1=0$, exactness gives
$W_0d_1=d_3^*V_0$. The top differential $d_3^*$ has trivial
kernel on analytic four-form coefficients, so $V_0d_0=0$.
A right lift gives $V_0=X_0d_1$. Thus
$W_1=W_0-d_3^*X_0$ satisfies both $W_1d_1=0$ and
$d_2^*W_1=E$. Lift once more to write $W_1=C_0d_2$,
and replace $C_0$ by $(C_0+C_0^*)/2$.
All these lifts act on the infinite lattice coefficients before
periodization; finite-period harmonic modes do not supply the
analytic exactness used here.

Put $B_c=H_c^{-1}$, $B_f=H_f^{-1}$ for the unmodified metrics,
and
\[
 Z_0=B_c-T^*B_fT-b^{-1}d_{1,c}d_{1,c}^*.
\]
Equation \eqref{ch05:physical-transports} gives
$Z_0H_cd_{1,c}=0$.
Applying the two-sided lift to $H_cZ_0H_c$ and the one-sided
lift to $d_{2,f}B_fTH_c$ produces
\begin{equation}
 Z_0=(d_{2,c}B_c)^*C_0(d_{2,c}B_c),\qquad
 Sd_{2,c}B_c=d_{2,f}B_fT.
 \label{ch05:defect-lifts}
\end{equation}
The estimates are
$\|C_0\|\le C_0'$ and $\|S^*S\|\le C_1L^{-2}$.
Indeed a fixed fine-residue entry of $T$ or $B_fT$ has
weighted coefficient bound $CL^{-2}$ and one fine difference
has bound $CL^{-3}$. A Gram product sums over $L^4$ residues.
Its bounds are thus $L^4(CL^{-2})^2=O(1)$ and
$L^4(CL^{-3})^2=O(L^{-2})$, respectively.
The lift acts on coarse exterior indices separately for every
fine residue, so it introduces no further residue sum.

Modify the inverse metrics as in
\eqref{ch05:metric-modification}, with coefficients
$\zeta_c,\zeta_f$. Write the modified inverse metrics as
$\widetilde B_c,\widetilde B_f$, and
$\widetilde H_c=\widetilde B_c^{-1}$.
Set
\[
 D_c=(I+\zeta_c d_{2,c}B_cd_{2,c}^*)^{-1},\qquad
 R_0=C_0+\zeta_c I-\zeta_f S^*S.
\]
Expansion of \eqref{ch05:defect-lifts} gives the exact identity
\begin{equation}
 \widetilde H_c
   (\widetilde B_c-T^*\widetilde B_fT
                         -b^{-1}d_{1,c}d_{1,c}^*)
       \widetilde H_c
       =d_{2,c}^*D_c^*R_0D_cd_{2,c}.
 \label{ch05:positive-defect}
\end{equation}
If $\zeta_c\ge C_0'+2$ and $L^2\ge\zeta_f C_1$, then
$R_0\ge I$. Since $D_c^{-1}$ has a uniform upper bound for
fixed $\zeta_c$, the three-form metric $D_c^*R_0D_c$ is
uniformly strictly positive. Its inverse and positive square
root have exponential coefficient bounds on a smaller strip
by the usual analytic functional calculus.
This proves positivity on the whole three-form space.

For each adjacent pair of stages choose the coarse modification
large enough for its finitely many incoming possibilities.
Fix these coefficients before choosing the block ratios.
Then the condition $L^2\ge\zeta_f C_1$ is a finite lower
bound on each permitted ratio. No constant uniform over all
smoothing orders is asserted: the estimates are uniform over
histories ending at a fixed order.

\section{Convergence at a fixed scale}
\label{ch05:horizon-convergence}

The curvature metrics and transport operators in this section are those of Sections~\ref{ch05:curvature} and~\ref{ch05:transports}; their histories now vary while the observation scale stays fixed.

The preceding uniform bounds also give convergence for the particular
nested histories used in the construction. The nesting is essential:
fix the ratios, smoothing orders and modification coefficients at every
scale, and increase only the terminal cutoff. At a fixed scale $m$, the
history with terminal scale $N+1$ extends the history with terminal scale
$N$ at its fine end. An arbitrary sequence of bounded histories need
not converge.

We first record how to identify limits on the infinite lattice. Write a
translation-invariant row as $K_{ab}(z)$, where $z\in\ZZ^4$ and $a,b$
range over fixed finite coordinate sets. For an adjacent-scale row,
include the finitely many fine residues in these coordinate sets. Suppose
\[
 \sum_z e^{\sigma|z|_1}|K_{N,ab}(z)|\le C
 \qquad(\sigma>0)
\]
uniformly in $N$. Every subsequence has a further subsequence converging
in each coefficient, by diagonal extraction. Its limit has the same
bound, by finite partial sums. For every $0\le\tau<\sigma$, coefficient
convergence implies
\begin{equation}
 \sum_z e^{\tau|z|_1}|K_{N,ab}(z)-K_{ab}(z)|\longrightarrow0.
 \label{ch05:horizon-weighted}
\end{equation}
Indeed each summand is bounded by
$2C e^{-(\sigma-\tau)|z|_1}$, a summable sequence. If the symbols have
a specified limit on a dense set of real momenta, any two cluster
symbols agree there, hence everywhere by continuity. Fourier inversion
then identifies every coefficient. Alternatively, agreement of cluster
matrices on an unbounded sequence of torus sides suffices: for fixed $z$,
the periodized entry is $\sum_h K(z+nh)$, and its terms with $h\ne0$
tend to zero as $n\to\infty$. Either criterion proves uniqueness of
the cluster limit and therefore convergence of the full sequence.
This argument identifies the lattice-cover coefficients, not merely a
matrix on one fixed torus.

\begin{proposition}[Convergence of the selected Gaussian rows]
\label{ch05:horizon-limit}
Fix a scale and the nested schedule just described. The unmodified
curvature metrics, their inverses, the adjacent one- and two-form
transports, and the modified metrics of
\eqref{ch05:metric-modification} converge in the sense of
\eqref{ch05:horizon-weighted}, at every width strictly smaller than a
common positive width available at that scale. The same conclusion
holds for the specified coordinate exterior lifts, their amended
three-form metric in \eqref{ch05:positive-defect}, and their positive
square roots. It also holds for a positive square root of
$\widetilde H-cI$ whenever the fixed direct coefficient $c$ is chosen
strictly below the uniform lower bound of $\widetilde H$.
All limits are for the actual chosen rows on the cover and their
periodizations. No choice of a subsequence, matrix eigenbasis or
square-root signs is made.
\end{proposition}

\begin{proof}
Let $d=N-m$ be the history depth and $M_{d,m}$ its total dilation.
Let $u_{d,m,\mu}(p)$ be the squared modulus of its accumulated directional
filter, evaluated at $p/M_{d,m}$. Thus $u_{0,m,\mu}=1$; the filter is the
product of the centered-box and link factors used in
\eqref{ch05:history-filter}. Write $g_{m,\mu}(\theta)$ for the
single-step directional filter at scale $m$, including the link average
$r_{L_m}(\theta_\mu)$. Composition gives
\[
 u_{d+1,m,\mu}(p)=|g_{m,\mu}(p/L_m)|^2
                         u_{d,m+1,\mu}(p/L_m).
\]
Every factor is an average of unit complex numbers on real momenta,
possibly multiplied by a unit phase. Its modulus is at most one.
Induction, with the schedule fixed at each scale, gives
$0\le u_{d+1,m,\mu}(p)\le u_{d,m,\mu}(p)\le1$.
Thus the directional filter squares converge at every real $p$, including
any fixed reciprocal translate.

The accumulated noise is also convergent. Extending the history adds one
independent noise contribution propagated through all the previously
fixed coarser steps. The old contributions are unchanged, and every
new contribution is nonnegative at real momenta. The accumulated noise
therefore increases with depth. The suffix estimate in
Proposition~\ref{ch05:history-prop} bounds it uniformly, and the noise
at the last blocking step contributes one. Its limit is finite and at
least one. This argument concerns the actual nested noise list; it is
not a monotonicity assertion for arbitrary histories of the same length.

For the initial covariance contribution, extend the finite nonprincipal
alias sum to all $z\in\ZZ^4\setminus\{0\}$ by zero outside the centered
representatives modulo $M=M_{d,m}$. Its summand at an eventually included
fixed $z$ is
\[
 \frac{u_{d,m,\mu}(p+2\pi z)}{M^2\lambda((p+2\pi z)/M)}.
\]
The denominator tends to $|p+2\pi z|^2$, which is positive for nonzero
$z$ and $p$ in the base cube. The alias estimates above dominate these
summands by $C\prod_\nu(1+|z_\nu|)^{-2}$, independently of depth.
Dominated convergence proves convergence of the sum. At $p\ne0$ the
principal summand converges by the same argument without an alias sum.
Consequently each directional covariance $c_{N,m,\mu}(p)$ has a finite
limit at real $p\ne0$, bounded below by $b^{-1}$.

The curvature metric is the rational expression
\eqref{ch05:H0} in these covariances and the exterior derivatives. Its
scalar quotient denominator is
$\sum_\mu|q_\mu(p)|^2/c_{N,m,\mu}(p)$; its limit is positive away from
zero. Thus its symbol converges there. The uniform analytic bounds
already proved and the cluster-limit criterion identify one limit on
the cover, including at zero. Inverses have the same conclusion.
This avoids assigning a finite massless covariance to the zero mode.

For adjacent transports the ratio $L=L_m$ is fixed as $N$ increases.
At $p\ne0$, every one of its finitely many aliases
$\theta_l=(p+2\pi l)/L$ is a nonzero real momentum. All covariance
factors in \eqref{ch05:transport-formula} therefore converge.
The denominators $\rho,\rho_l$ have positive limits, and so does
\[
 \alpha=\sum_l |m_l|^2\tr D_l:
\]
the principal centered-box factor is nonzero on the real base cube and
$D_l$ has a positive diagonal limit. Hence the finite alias formulas
and their Fourier reconstruction converge. The previously proved
uniform exponential bounds identify unique cover limits for $W,T$.
At exceptional zero momenta uniqueness follows from continuity of
cluster symbols, rather than division by a vanishing expression.

Products preserve this convergence by the weighted convolution
inequality, with the fixed finite coordinate multiplicity included.
For uniformly bounded inverse rows, cluster limits satisfy the inverse
identity on each torus; uniqueness of the finite-dimensional inverse
and recovery from unbounded periods identify their cover limit.
The coordinate exterior lifts used in \eqref{ch05:defect-lifts} are
also continuous after a width decrease. Each elementary operation is
coordinate evaluation or a divided difference
\[
 \frac{\widehat K_N(p)-\widehat K_N(p|_{p_a=0})}{e^{ip_a}-1}.
\]
Its numerator converges, its denominator is nonzero on a dense set,
and its smaller-width bound was proved in the exterior construction.
The cluster criterion consequently applies through the finitely many
coordinate operations. This requires the fixed coordinate lift of that
construction, not a fresh arbitrary solution of the exterior identity.

It follows that $C_0,S,D_c$ and $D_c^*R_0D_c$ in
\eqref{ch05:positive-defect} converge, as do the modified metrics and
inverses. Finally consider one of the stated positive square roots.
Its uniform analytic bound gives cluster rows. On each finite torus
their matrices are positive semidefinite and have the same square,
namely the limit of the underlying positive matrix. The positive
semidefinite square root is unique by the finite-dimensional spectral
theorem. Thus the cluster matrices agree on every torus, and period
recovery identifies the cover row. This proves convergence for the
chosen roots themselves, not merely for their squares.
\end{proof}

In particular all canonical path-row coefficients obtained from these
rows converge, with a strictly positive gap between the available width and
the width used to sum paths. Finite products, fixed exterior operations
and the fixed positive factorizations preserve this conclusion. The
actual block average correction introduced below is a fixed finite-range
row at the selected scale, so its application to the convergent
straight transport also preserves convergence. These are statements
about the quadratic data. They do not imply cutoff continuity of the
nonlinear activity arrays or of marked integral maps; those require
their own stability arguments.

\section{Uniform inverse bounds: the hypotheses that matter}
\label{ch05:inverse}

A spectral gap in Euclidean norm is not, by itself, a bound for an
inverse in supremum norm independent of volume. One sufficient
framework is an exponentially local matrix with a uniform positive
gap. To state its quantitative content, assign each coordinate $i$
a lattice site $a_i$, with uniformly bounded coordinate multiplicity
per site, and define
\[
 \|T\|_\eta=\max\left\{
 \sup_i\sum_j e^{\eta d(a_i,a_j)}|T_{ij}|,
 \sup_j\sum_i e^{\eta d(a_i,a_j)}|T_{ij}|\right\}.
\]
Here $d$ is periodic lattice distance. Suppose $H=H^*\ge mI$ and
$\|H\|_{\eta_0}\le C$ uniformly in the volume.
For some $0<\eta<\eta_0$, depending only on $m,C,\eta_0$ and the
coordinate multiplicity, $H^{-1}$ has uniformly bounded
exponentially weighted row and column sums, after reducing the
exponent once more if necessary. This is the usual weighted
resolvent estimate behind exponentially local lattice propagators;
compare \cite{Balaban1984Propagators}.
It follows by conjugating with
$e^{\eta d(a_i,a_{i_0})}$ and applying the resolvent identity to the
resulting small perturbation. In particular
$\|H^{-1}\|_{\infty\to\infty}$ is bounded independently of volume.

For analytic momentum-dependent matrices, an analogous perturbation
argument must use holomorphic coordinates. Suppose $K(p)$ is
holomorphic in a strip of width $\sigma$ and, in fixed coordinates,
$\|K(p)\|\le A$ there, while at real $x$ its least singular value
is at least $a>0$. Cauchy estimates and the inverse perturbation
lemma give a smaller strip, depending only on $a,A,\sigma$,
on which $K^{-1}$ is holomorphic and uniformly bounded. This
statement is dimension independent.

In weighted Fourier spaces the weights often involve
$|q(\operatorname{Re}p)|$ and are not holomorphic. The preceding
argument is justified if, at every real anchor $x$, the matrices
defining the weighted norms have holomorphic replacements in a
fixed neighborhood of $x$, uniformly equivalent to the original
norms, and coinciding with them on the vertical slice
$\operatorname{Re}p=x$. One applies Cauchy's estimate in these
replacement coordinates. A bound in nonholomorphic coordinates
alone does not justify differentiating the weighted matrix.
This distinction is essential when the weight degenerates with a
scale parameter.

\section{The actual block average and the forest coordinates}
\label{ch05:actual-coordinates}

We now identify the nonlinear coordinates behind the Gaussian calculation.
Fix $j,L$, a coarse period $s\ge100(j+8)$, and in each centered block
$Ly+\mathbb B_{(L-1)/2}$ a spanning tree rooted at $Ly$. Trees contain
only internal block edges. Denote by $\iota_F$ extension from free-edge
coordinates to all fine edges, with value zero on tree edges. For coarse
links $V$, their replica is
\[
 (\operatorname{rep}V)_e=
 \begin{cases}1,&\operatorname{parent}(e_-)=\operatorname{parent}(e_+),\\
 V_{y,\mu},&\operatorname{parent}(e_-)=y,
                    \operatorname{parent}(e_+)=y+e_\mu.
 \end{cases}
\]
The centered-block convention and odd $L$ make these the only possibilities
for a positive fine edge in direction $\mu$. In particular the replica
is one on every tree edge. The actual forest chart is
\begin{equation}
 U(v,\psi)_e=\exp((\iota_F\psi)_e)
                         (\operatorname{rep}\exp v)_e.
 \label{ch05:forest-chart}
\end{equation}
It is a chart on the tree slice, not a choice of independent variables on
tree edges. Its derivative at the identity is exactly
\begin{equation}
 DU(0,0)[h,z]=\operatorname{rep}h+\iota_Fz,
 \qquad DU(0,0)[0,z]=\iota_Fz.
 \label{ch05:forest-derivative}
\end{equation}
All color coordinates are orthonormal for the trace metric. Thus
$\|\iota_Fz\|_2=\|z\|_2$ in counting norms. This identifies the vertical
coordinates used by the Hessian with the arrays in the tree inequality
below; no volume-dependent coordinate conversion occurs.

\subsection{A common analytic branch for all local labels}

The minimizing frame and block average of Chapter~\ref{ch04:blocking} are
real analytic near the identity. For the frame this follows directly from
its chart equation: its quadratic objective at zero is
$\tfrac12\operatorname{av}_{r,r'}|h_r-h_{r'}|^2$, positive definite after
$h_0=0$ is imposed. For the block average, its Hessian in the group variable
is the identity. Complexify these finite-dimensional real analytic chart
equations and apply the analytic implicit-function theorem
\cite[Theorem 1.3.8]{LeblSCV}. The resulting critical points are the unique
minimizers already identified in Lemma~\ref{ch04:frame-lemma}, after
reducing the neighborhood if necessary.

At fixed $j,L$ every such equation reads a translate of one fixed finite
patch. Choose a radius in the supremum norm of all its input logarithms
on which these analytic extensions and their derivatives are bounded.
There are finitely many directions and local tree shapes at fixed $L$;
take their common radius and bounds. Periodic pullback to a patch has
supremum norm at most one, so these bounds are independent of $s$.
Decrease this radius until every loop test and minimum-cost test is
strictly on its minimizing branch. Every local label then selects the
same frame and block average as a function of the full fine field on the
whole common ball. Redundant uniform branch masses are constant there.
This is equality of the local action functions before stationarity is
imposed, not merely equality along a stationary graph. It concerns the
uncut action; shrinking smooth cutoff plateaux are treated separately.

Let $a$ be an infinitesimal fine field. Write $I_a(r\to r';o)$ for its
signed sum along a coordinate path in the frame patch at $Ly$. The
linearized minimizing-frame equation gives
\begin{align}
 \tau_y(r)&=-\operatorname{av}_{r',o}I_a(Ly+r\to Ly+r';o),\notag\\
 H_{y,r}(a)&=\tau_y(r)-\tau_y(0),&
 R a(y)&=-\sum_rw_r^{\rm av}H_{y,r}(a).
 \label{ch05:frame-linearization}
\end{align}
Indeed its equation is $H_r-\operatorname{av}_{r'}H_{r'}
=-\operatorname{av}_{r',o}I_a(r\to r';o)$, and $H_0=0$ fixes the additive
constant. Differentiating the actual connector and then its block average
therefore yields
\begin{equation}
 \mathcal Q a=M a+d_{0,c}R a,
 \qquad \mathcal Q d_{0,f}\phi=d_{0,c}(\phi|_{L\Lambda_s}).
 \label{ch05:actual-Q}
\end{equation}
Here $M$ is exactly the straight length-$L$ average with the convolution
weights used in \eqref{ch05:alias-recursion}. The second identity follows
also by differentiating gauge covariance at block roots. Each route has
length $O_j(L)$; hence $R$ has finite coarse range and absolute row sum
at most $C_jL$. These statements include the root normalization. Dropping
it would change $\mathcal Q$.

\subsection{The tree comparison and the actual Hessian}

For a fine edge $e$, join its endpoints to their block roots along the
chosen trees. In the tree slice the resulting walk has transport $U_e$.
It lies in one or two adjacent blocks and has length at most $2L^4-1$.
In the macro axial gauge of Section~\ref{ch04:predictor-geometry}, its
distance from the identity is bounded by this length times the sum of bond distances from the identity in the macro-box. The
first-power frame estimate controls the two root-to-frame discrepancies;
the block average estimate controls the coarse root-to-root transport. Thus
\begin{equation}
 \chi\bigl(U_e(\operatorname{rep}V)_e^{-1}\bigr)
 \le C_{j,L}\mathcal P_y(U)
       +\mathbf1_{\{e\text{ crosses blocks}\}}
                            \chi(V_{y,\mu}Q_{y,\mu}(U)^{-1}).
 \label{ch05:nonlinear-tree}
\end{equation}
For an internal edge the walk closes at its root and only the first term
is needed. For a crossing edge, insert the block average between the two
root frames and use the triangle inequality for the normalized Hilbert--Schmidt distance. The endpoint errors
are bounded by the sum in \eqref{ch04:frame-first-power}, so this step
requires no pointwise frame estimate not already proved. The constant is
allowed to depend on $L$; the stronger pre-$L$ factor in
\eqref{ch04:macro-comparison} is not required for this use.

Apply \eqref{ch05:nonlinear-tree} to $U=\exp(\epsilon a)$ with $a$ zero
on tree edges, and $V=1$. Divide by $|\epsilon|$ and let $\epsilon\to0$.
The first-order variation of the distance from the identity is the color norm, so
\[
 |a_e|\le C_{j,L}\sum_{p\text{ in the macro-patch at }y}|(d_{1,f}a)_p|
       +\mathbf1_{\{e\text{ crosses blocks}\}}|(\mathcal Qa)_{y,\mu}|.
\]
Square, use the finite macro-patch cardinality, and sum over edges. A fine
plaquette belongs to a bounded number of macro-patches, and each coarse
edge occurs for at most a fixed number of fine edges at fixed $L$.
Consequently
\begin{equation}
 \|a\|_2^2\le B_{j,L}\|d_{1,f}a\|_2^2
                       +C_{j,L}\|\mathcal Qa\|_2^2,
 \quad a=\iota_Fz.
 \label{ch05:tree-comparison}
\end{equation}
Both constants are uniform in the period and in the allowed trees.
The use of periodic lifts preserves the bounded multiplicity.

Let $\mathcal A_f$ be the uncut fine action, without the common inverse
coupling or inherited loop action, and set
\begin{equation}
 f(v,\psi)=\mathcal A_f(U(v,\psi))
       +b\sum_{e\text{ coarse}}w\bigl((\exp v)_e
                                   Q_e(U(v,\psi))^{-1}\bigr).
 \label{ch05:actual-action}
\end{equation}
Assume its actual quadratic Taylor coefficient is
$K_f=d_{1,f}^*H_fd_{1,f}$ with $H_f\ge\varepsilon I>0$. This coefficient
is read from the same path rows in $\mathcal A_f$: each quadratic path
square contributes the square of its linearized curvature row, and atoms
of curvature degree at least three contribute nothing at second order.
Identifying their sum with a chosen metric $H_f$ is an exact coefficient
condition on the incoming path action.

The first derivatives of both terms of \eqref{ch05:actual-action} vanish
at the identity. Therefore the second derivative of the forest chart does
not enter its Hessian. By \eqref{ch05:forest-derivative}, for
$a=\operatorname{rep}h+\iota_Fz$ and
$a'=\operatorname{rep}h'+\iota_Fz'$, it is precisely
\begin{equation}
 D^2 f(0)[(h,z),(h',z')]
   =\langle d_{1,f}a,H_fd_{1,f}a'\rangle
           +b\langle h-\mathcal Qa,h'-\mathcal Qa'\rangle.
 \label{ch05:actual-hessian}
\end{equation}
No derivative of the nonlinear block average beyond its actual first
derivative is missing here: its residual vanishes at zero. In particular,
\begin{equation}
 D_\psi^2f(0,0)\ge
       \frac{\min(\varepsilon,b)}{1+B_{j,L}+C_{j,L}}I.
 \label{ch05:actual-gap}
\end{equation}
This proves the vertical gap for the actual forest chart, rather than
assuming it from the flat Fourier calculation.

\subsection{A smooth representative of the stationary tangent}
\label{ch05:actual-tangent}

Let $W_s$ denote the straight-average transport already constructed in
Section~\ref{ch05:transports}. Its identities in physical coordinates are
\[
 K_fW_s=M^*K_c,\qquad MW_s=I-b^{-1}K_c.
\]
These hold for the same incoming metric and penalty as
\eqref{ch05:actual-hessian}. Since $K_cd_{0,c}=0$, replacing $M$ by
$\mathcal Q=M+d_{0,c}R$ does not change the right side of the first
identity. The second identity does require a gradient correction.

Choose a smooth compactly supported cardinal partition of unity
$\varphi$ on $\RR^4$: $\varphi(n)=\mathbf1_{n=0}$ for $n\in\ZZ^4$ and
$\sum_{n\in\ZZ^4}\varphi(x-n)=1$. For example, in one dimension divide a
smooth bump supported in $(-1,1)$ and positive there by the sum of its
integer translates, and take the fourfold product. Define a periodic
extension of coarse site fields to fine sites by
\[
 (E h)(x)=\sum_{y\in\Lambda_s}\sum_{k\in\ZZ^4}
                \varphi(x/L-y-sk)h(y).
\]
It equals $h(y)$ at $x=Ly$, is supported within one coarse block of its
input, and has fine differences bounded by
$\|\nabla_f^rE\|_{\infty\to\infty}\le C_rL^{-r}$. Now put
\begin{equation}
 W_a=W_s-d_{0,f}ERW_s.
 \label{ch05:corrected-transport}
\end{equation}
By \eqref{ch05:actual-Q},
\begin{equation}
 \mathcal QW_a=MW_s=I-b^{-1}K_c,
 \qquad K_fW_a=\mathcal Q^*K_c,
 \qquad d_{1,f}W_a=d_{1,f}W_s.
 \label{ch05:actual-transport-identities}
\end{equation}
Thus $W_a$ minimizes the actual quadratic energy and has the same reduced
quadratic form $K_c$. This holds also on constant modes: $E$ extends
constants to constants, gradients have zero curl, and no inverse of a
zero-momentum Laplacian has been introduced.

The correction preserves the needed smooth interpolation bounds. The
finite-range row $R$ has size $C_jL$ and the already proved kernel of
$W_s$ has size $C_jL^{-1}$ with exponential coarse decay. Their composition
$RW_s$ therefore has an $O_j(1)$ exponentially decaying coarse kernel.
Applying $r+1$ fine differences to $E$ gives
\begin{equation}
 |\nabla_f^rW_a(Ly+x,y')|
       \le C_{j,r}L^{-1-r}e^{-c_jd(y,y')},
 \quad r\le j-2,
 \label{ch05:actual-bridge}
\end{equation}
with the same permitted orders as for $W_s$. The extension adds only
bounded coarse range, so it preserves exponential summability. This
argument uses finite-difference bounds of a cardinal extension, not an
assumption that a discontinuous block-constant extension is smooth.

Finally adjust $W_av$ by the unique gradient $d_{0,f}\phi$ whose site
potential is zero at every block root and makes all tree edges vanish.
It is obtained successively along each tree. Such a gradient leaves
$\mathcal QW_av$ unchanged by \eqref{ch05:actual-Q}, and leaves its curl
unchanged. Write the resulting representative as $W_Fv$. Coercivity
\eqref{ch05:actual-gap} makes it the unique minimizer in the forest slice.
Once a stationary graph $\psi=\rho(v)$ has been constructed, differentiating
its equation and using \eqref{ch05:actual-hessian} proves
\begin{equation}
 \operatorname{rep}v+\iota_FD\rho(0)v=W_Fv.
 \label{ch05:stationary-tangent}
\end{equation}
The smooth representative $W_a$ and the forest representative $W_F$ are
distinct. Estimate~\eqref{ch05:actual-bridge} is asserted for $W_a$, not
for tree coordinates. Gauge-invariant quadratic insertions see the same
curl in either representative.

\section{A stationary background on a uniform neighborhood}
\label{ch05:stationary}

We first state the quantitative implicit-function estimate needed
for the nonlinear action. Let $x$ be the coarse coordinates and
$y$ the free fine coordinates after gauge fixing. Both carry the
supremum norm. Write $f(x,y)$ for the action including the block
penalty, and $F(x,y)=D_yf(x,y)$.
All norms below are operator norms induced by these supremum norms.

\begin{proposition}[Uniform stationary graph]
\label{ch05:stationary-prop}
Suppose $F$ is holomorphic on a product neighborhood
$\|x\|_\infty<r_0$, $\|y\|_\infty<r_0$, with $F(0,0)=0$.
Let $H_*=D_yF(0,0)$ be invertible and assume, uniformly in the
number of coordinates,
\[
 \|H_*^{-1}\|\le J,\qquad
 \|D_yF(x,y)-H_*\|\le C\max(\|x\|_\infty,\|y\|_\infty),
 \qquad \|D_xF(x,y)\|\le B .
\]
Choose $0<r_y<r_0$ and $0<r_x\le r_y$ so that
$JC r_y\le1/4$ and $JB r_x\le r_y/4$.
Then for $\|x\|_\infty<r_x$ there is a unique solution
$y=\rho(x)$ of $F(x,y)=0$ in $\|y\|_\infty\le r_y$.
It is holomorphic, satisfies $\rho(0)=0$ and
$\|\rho(x)\|_\infty\le r_y/2$, and obeys
\[
 D\rho(x)=-[D_yF(x,\rho(x))]^{-1}D_xF(x,\rho(x)),
 \qquad \|D\rho(x)\|\le \tfrac43 JB .
\]
\end{proposition}

This is the quantitative contraction form of the implicit-function
theorem: apply it to $y-H_*^{-1}F(x,y)$ on the closed $r_y$-ball.
The two displayed smallness conditions give contraction constant
$1/4$ and displacement of the origin at most $r_y/4$.
We use the standard theorem rather than repeat its proof; holomorphic regularity is the analytic implicit-function theorem \cite[Theorem 1.3.8]{LeblSCV}. The real
smooth version is also a direct application of the implicit-function
theorem \cite[Theorem 9.28]{Rudin1976}.
For a merely $C^N$ real action the conclusion is $C^{N-1}$ for
$\rho$, not holomorphy.

In applications the operator bounds in this proposition should be
deduced from local derivative estimates before any dimension
dependent conversion of norms. For example, suppose each coarse or fine coordinate
has an anchor $a_i$, the total number of coordinates anchored at any
site is bounded uniformly in the volume, and, for $2\le q\le N$, every rooted derivative
satisfies
\[
 |\partial_{i_1}\cdots\partial_{i_q}f(w)|
 \le C\exp\left(-\eta\sum_{h=1}^q d(a_{i_t},a_{i_h})\right)
 \quad\text{for every }t.
\]
The lattice shell sum
$S_\eta=\sup_v\sum_i e^{-\eta d(v,a_i)}$ is bounded uniformly in
the torus size. Summing the second derivative gives $B\le CS_\eta$;
summing the third gives the Hessian Lipschitz bound
$CS_\eta^2$. A real Hessian gap together with exponential locality
supplies the required $J$ by the previous section.
These implications yield a uniform real neighborhood. A uniform
complex neighborhood requires the corresponding complex estimates,
which do not follow from a fixed finite number of real derivatives.

At fixed derivative order, higher derivatives of the background and
of its Hessian can then be controlled by repeated implicit
differentiation. If only polynomial dependence on an auxiliary
volume is needed, this conversion can be postponed until after the
uniform neighborhood and the first derivative have been obtained.
It must not be used to choose the neighborhood itself.
For a positive real Hessian $H(x)$, the determinant contribution is
differentiated by the standard identity
$D\log\det H[h]=\tr(H^{-1}DH[h])$.

\subsection{Supplying the local derivative estimates}

We retain the distinction from Section~\ref{ch05:actual-tangent}: $W_a$ is the smooth representative of the actual stationary tangent, whereas $W_F$ is its representative in the forest coordinates. They have the same curvature, but the fine-difference bounds apply to $W_a$.

The derivative hypotheses have a direct origin in the
exponentially weighted path expansion. We spell it out because
a bound for the action on each finite torus would not supply a
uniform neighborhood.

Suppose the uncut classical action is a sum of local analytic
terms indexed by rooted finite path collections. Let $\ell$
be the total path length plus one. Require that each collection
have connected support containing its root, and that
\begin{equation}
 \sup_z\sum_{\gamma:\ {\rm root}(\gamma)=z}
                   |c_\gamma|e^{\sigma\ell(\gamma)}\le P .
 \label{ch05:path-budget}
\end{equation}
A bounded number of paths per term is sufficient; products of
their matrix entries, normalized traces and fixed invariant
contractions are allowed. The constants $\sigma,P$ and this
number are independent of the torus.
Choose local coordinate maps with uniformly bounded analytic
extensions and fixed support radius at the chosen block ratio.

A complex coordinate of modulus at most $r$ changes each matrix
factor's norm by at most $e^{Cr}$. Therefore a term of total
length $\ell$ has complex modulus at most $Ce^{Cr\ell}$.
Choose $r$ so that $Cr<\sigma/4$. Differentiating $q$ times
selects at most $C_q(1+\ell)^q$ positions; the same estimates
hold on the smaller complex ball. If the differentiated
coordinates have prescribed anchors $a_{i_1},\ldots,a_{i_q}$,
connectedness implies
\[
 \sum_{h=1}^q d(a_{i_t},a_{i_h})\le C_R q(1+\ell)
             \quad\text{for every }t .
\]
Here $C_R$ allows for the finite support radius of the coordinate
maps, fixed at this block ratio. A fraction of the exponential
weight in \eqref{ch05:path-budget} controls this distance, and a
further fraction absorbs the polynomial $(1+\ell)^q$.
The remaining exponential weight bounds the sum over possible
roots of a term visiting a fixed anchor: the number of roots
within distance $C_R(1+\ell)$ grows only polynomially in $\ell$.
It follows, for every fixed $N$, that
\[
 |\partial_{i_1}\cdots\partial_{i_q}f(w)|
 \le C_N\exp\left(-\eta_N
                 \sum_h d(a_{i_t},a_{i_h})\right),
 \qquad 2\le q\le N ,
\]
on a common complex ball. Block penalties and
local coordinate terms add finite-range contributions satisfying
the same estimate. Smooth cutoffs are not being analytically
continued: the action here is the uncut action whose germ at
zero defines the background.

For the actual forest chart \eqref{ch05:forest-chart}, the deterministic
tree comparison and the Hessian identification have already been proved
in \eqref{ch05:tree-comparison}--\eqref{ch05:actual-gap}. The chosen
block average is analytic on a common fixed-patch ball, and all its redundant
local labels agree there as functions of the full field.
The derivative bounds above give exponential locality of that
Hessian. Section~\ref{ch05:inverse} now supplies a supremum-norm
bound for its inverse independent of volume; a sufficiently
small complex perturbation preserves this bound.
Proposition~\ref{ch05:stationary-prop} therefore applies on a
uniform neighborhood.

This argument proves the uniform stationary construction for
uncut path actions satisfying \eqref{ch05:path-budget}, the
actual incoming quadratic coefficient condition in
\eqref{ch05:actual-hessian}. The actual block average and forest coordinates
supply the local analytic-coordinate hypotheses and the deterministic
comparison \eqref{ch05:tree-comparison}. It does not require
a later matching theorem. If two local descriptions agree
as functions on the same coordinate neighborhood, uniqueness
makes their stationary graphs agree there.
Equation~\eqref{ch05:stationary-tangent} identifies the tangent with
the forest representative of the actual corrected transport $W_a$.
The stationary chain rule and \eqref{ch05:actual-transport-identities}
give the coarse quadratic form as the Hessian of the reduced action. Equality of actions only on a
previously selected graph would not suffice for this argument.

The constants are uniform in the period after the block ratio,
smoothing order and path bound have been fixed. Uniformity in
an arbitrary block ratio is a separate weighted-norm estimate;
the Euclidean tree gap above is allowed to depend on $L$.

\begin{proposition}[Uniform actual stationary chart]
\label{ch05:actual-stationary}
Fix the group, $j,L,b$, an exponential path bound
\eqref{ch05:path-budget}, and the actual incoming quadratic metric with
$H_f\ge\varepsilon I>0$. For the uncut action
\eqref{ch05:actual-action} there are radii $r_V,r_\psi>0$ and constants
$c,J,C_\rho>0$, independent of the allowed period and block trees, such
that its unique stationary graph $\rho$ satisfies
\[
 \|\rho(v)\|_\infty<r_\psi/2,
 \qquad \|D\rho(v)\|_{\infty\to\infty}\le C_\rho
 \quad(\|v\|_\infty<r_V).
\]
On the corresponding smaller real product chart the vertical Hessian is
at least $cI$ and its inverse has supremum operator norm at most $J$.
All local tag descriptions define the same action and stationary graph
there. For every prescribed finite order $N$, ordinary derivatives of
$f$ through order $N$, of $\rho$ through order $N-1$, and of the centered
log determinant of its vertical Hessian through order $N-2$ are bounded
by fixed polynomials in the total coordinate number $n$. The radii and
first-derivative bound do not acquire a factor depending on $n$.
\end{proposition}

To check the last assertions, the third rooted derivative bound gives
$\|H(w)-H(0)\|_{2\to2}\le C\|w\|_\infty$ by row and column sums.
Reduce the common ball until this is at most half the gap in
\eqref{ch05:actual-gap}. The weighted inverse argument then applies
uniformly throughout that real ball. Proposition~\ref{ch05:stationary-prop}
gives the graph and its volume-independent Lipschitz constant. Repeated
implicit differentiation isolates $H D^k\rho$ at order $k$; every other
term involves lower derivatives of $\rho$ and derivatives of $f$ of order
at most $k+1$. Converting bounded coordinate entries to ordinary
multilinear norms costs only a fixed power of $n$ at fixed order. Induction
therefore proves the asserted polynomial bounds. The identities
$D H^{-1}=-H^{-1}(DH)H^{-1}$ and
$D\log\det H=\tr(H^{-1}DH)$ give the determinant assertion, with one
additional factor $n$ for the trace. These are standard implicit and
Gaussian determinant differentiation rules; no polynomial-in-volume
bound is used to select the radii. For an application requiring $k+4$
derivatives, take the path/coordinate derivative order at least $k+6$
so that the stated determinant margin is retained.

\section{Scope of the completed Gaussian construction}
\label{ch05:pending}

For each fixed smoothing order, the finite-history construction
now provides uniformly analytic positive curvature metrics,
the specified one- and two-form transports, and a positive
three-form factor for the amended curvature defect. Its
parameter choices precede the choice of the scale sequence.
The nonlinear stationary graph is supplied for each incoming
path action with the explicit local bound and actual quadratic
coefficient condition just stated; its tree comparison and block average
branch compatibility are now proved for the chosen block kernel.

For the nested selected schedule, Proposition~\ref{ch05:horizon-limit}
now also proves fixed-scale convergence of the actual quadratic rows
at every smaller exponential width. For the actual nonlinear
step, the chosen path representation and block average must be
checked against \eqref{ch05:path-budget} and the actual quadratic
coefficient condition \eqref{ch05:actual-hessian}. The common minimizing
branch and equality of its local descriptions have been checked above. These are precise
interface conditions, not consequences of Gaussian algebra.
These conditions will be imposed on the admissible path action and
retained by the nonlinear step; a Gaussian identity alone would not
justify dropping them.
